\documentclass[12pt,a4paper]{article}

\usepackage{fontspec}
\usepackage{polyglossia}
\setdefaultlanguage{ukrainian}
\setmainfont{Helvetica}
\newfontfamily\cyrillicfont[Script=Cyrillic]{Arial}
\setmonofont{Menlo}[Scale=0.86]
\newfontfamily\cyrillicfonttt[Script=Cyrillic]{Arial}

\usepackage[a4paper,left=20mm,right=20mm,top=17mm,bottom=16mm]{geometry}
\usepackage{xcolor}
\usepackage{fancyhdr}
\usepackage{enumitem}
\usepackage{array}
\usepackage{tabularx}
\usepackage{ragged2e}
\usepackage{setspace}
\usepackage{hyperref}

\definecolor{accent}{HTML}{00636B}
\definecolor{codebg}{HTML}{EEF4F6}
\hypersetup{colorlinks=true,urlcolor=accent,linkcolor=accent}

\setlength{\parindent}{0pt}
\setlength{\parskip}{4pt}
\setstretch{1.04}
\setlist{leftmargin=6mm,itemsep=2pt,topsep=2pt,parsep=0pt}
\setlength{\headheight}{13pt}
\emergencystretch=2em

\pagestyle{fancy}
\fancyhf{}
\fancyhead[C]{\color{accent}\rule{\textwidth}{2.4pt}}
\fancyfoot[L]{\fontsize{7}{8}\selectfont ЛАБОРАТОРНА РОБОТА № 1 / МЕТОДИЧНІ ВКАЗІВКИ}
\fancyfoot[R]{\fontsize{7}{8}\selectfont\thepage}
\renewcommand{\headrulewidth}{0pt}
\renewcommand{\footrulewidth}{0pt}

\newcommand{\eyebrow}[1]{%
  {\color{accent}\fontsize{7.5}{9}\selectfont\bfseries\MakeUppercase{#1}\par}\vspace{5pt}}
\newcommand{\pagetitle}[1]{%
  {\fontsize{18}{21}\selectfont\bfseries #1\par}\vspace{7pt}}
\newcommand{\subhead}[1]{%
  \vspace{4pt}{\fontsize{10}{12}\selectfont\bfseries\MakeUppercase{#1}\par}\vspace{1pt}}
\newcommand{\smallhead}[1]{%
  \vspace{3pt}{\color{accent}\fontsize{8}{10}\selectfont\bfseries\MakeUppercase{#1}\par}\vspace{1pt}}
\newcommand{\cmd}[1]{%
  \par\vspace{2pt}\noindent\colorbox{codebg}{%
    \parbox{\dimexpr\linewidth-2\fboxsep\relax}{\ttfamily\small #1}}\par\vspace{2pt}}
\newcommand{\filename}[1]{\texttt{#1}}
\newcommand{\course}{Інженерія програмного забезпечення та мови програмування}

\begin{document}

\eyebrow{Методичні вказівки до самостійного виконання}
{\fontsize{18}{21}\selectfont\bfseries Лабораторна робота № 1\par}
\vspace{5pt}
{\fontsize{14}{17}\selectfont\bfseries Імперативна та структурна парадигми програмування\par}
\vspace{5pt}

Дисципліна «\course»

Спеціальність «Інженерія програмного забезпечення»

Перший курс

\subhead{Мета роботи}

Зрозуміти відмінність між неструктурованим імперативним керуванням і
структурним програмуванням; навчитися аналізувати переходи \texttt{goto},
відновлювати логіку алгоритму та замінювати переходи послідовністю,
розгалуженням і циклом без зміни поведінки програми.

\subhead{Результат роботи}

Структурована версія індивідуальної програми мовою C без операторів
\texttt{goto}; схема або таблиця початкових переходів; підтвердження однакового
виведення початкової та переробленої програм; оформлений звіт.

\subhead{Програмне забезпечення}

Windows, Linux або macOS; Visual Studio Code чи інший редактор; компілятор C
з підтримкою C11 (GCC або Clang). Початковий файл
\filename{task\_XX\_назва.c} надає викладач. Номер \texttt{XX} є номером
індивідуального варіанта.

Викладач: Петров О.В., к.т.н., доцент.

\subhead{Зміст}

\begin{enumerate}
  \item Підготовка та початковий запуск — с. 2.
  \item Обов’язкове завдання — с. 3.
  \item Перевірка та самоконтроль — с. 4.
  \item Структура й оформлення звіту — с. 5–6.
  \item Зразок титульного аркуша — с. 7.
  \item Здавання та критерії перевірки — с. 8.
  \item Установлення компілятора C у Windows — с. 9.
\end{enumerate}

\newpage

\eyebrow{Порядок виконання}
\pagetitle{1. Підготовка до роботи}

\smallhead{1.1. Отримайте свій варіант}

Викладач призначає кожному студентові один файл із діапазону
\filename{task\_01\_*.c}—\filename{task\_50\_*.c}. Працюйте лише зі своїм
варіантом. Не обмінюйте початковий файл на інший після початку роботи.

\smallhead{1.2. Створіть робочу папку}

Створіть папку \filename{Lab01\_Прізвище\_Група}. Скопіюйте до неї призначений
файл двічі та назвіть копії:

\begin{itemize}
  \item \filename{task\_XX\_original.c} — незмінний початковий код;
  \item \filename{task\_XX\_structured.c} — файл для перероблення.
\end{itemize}

Відкрийте всю папку у VS Code командою \textbf{File → Open Folder}.

\smallhead{1.3. Перевірте компілятор}

Відкрийте \textbf{Terminal → New Terminal} і виконайте одну з команд:

\cmd{gcc --version}
\cmd{clang --version}

Достатньо одного доступного компілятора. Якщо обидві команди не
розпізнаються, скористайтеся інструкцією для Windows на с. 9 або зверніться до
викладача щодо налаштування вашої операційної системи.

\smallhead{1.4. Скомпілюйте початкову програму}

Замініть \texttt{XX} своїм номером і виконайте:

\cmd{gcc -std=c11 -Wall -Wextra -Wpedantic task\_XX\_original.c -o original}

Якщо використовуєте Clang, замініть лише слово \texttt{gcc} на
\texttt{clang}. Компіляція має завершитися без помилок і попереджень.

\smallhead{1.5. Запустіть і збережіть еталонне виведення}

У Windows PowerShell виконайте:

\cmd{.\textbackslash original.exe | Tee-Object original\_output.txt}

У Linux або macOS виконайте:

\cmd{./original | tee original\_output.txt}

Прочитайте виведення та коротко сформулюйте призначення програми.
Надалі не змінюйте файл \filename{task\_XX\_original.c}.

\newpage

\eyebrow{Виконайте всі пункти}
\pagetitle{2. Обов’язкове завдання}

\smallhead{2.1. Проаналізуйте початкове керування}

Знайдіть усі мітки та оператори \texttt{goto}. Для кожної мітки визначте,
звідки до неї переходить виконання і за якої умови. Виділіть:

\begin{itemize}
  \item переходи вперед, що утворюють розгалуження;
  \item переходи назад, що утворюють цикли;
  \item переходи до завершення циклу або всієї програми;
  \item безумовні переходи та переходи всередині \texttt{if}.
\end{itemize}

Подайте аналіз як схему керування або таблицю «мітка — джерело переходу —
умова — роль у алгоритмі».

\smallhead{2.2. Сплануйте структурну заміну}

Для кожного фрагмента оберіть одну з базових структур:

\begin{itemize}
  \item \textbf{послідовність} — команди виконуються одна за одною;
  \item \textbf{розгалуження} — \texttt{if}, \texttt{if--else} або вкладені умови;
  \item \textbf{повторення} — \texttt{while}, \texttt{do--while} або \texttt{for};
  \item \texttt{break} чи \texttt{continue}, якщо вони безпосередньо передають
        зміст виходу з циклу або переходу до наступної ітерації.
\end{itemize}

\smallhead{2.3. Переробіть програму}

Редагуйте лише \filename{task\_XX\_structured.c}. Видаліть усі мітки та
оператори \texttt{goto}, замінивши їх структурними конструкціями. Збережіть:

\begin{itemize}
  \item початкові дані, числові константи та порядок їх опрацювання;
  \item формули, умови, накопичення значень і лічильники;
  \item тексти повідомлень та порядок виведення;
  \item усі особливі випадки: пропуск, відхилення, обмеження й завершення.
\end{itemize}

Можна вводити допоміжні логічні змінні та виділяти осмислені функції. Не
можна підставляти наперед обчислену відповідь, вилучати гілки або імітувати
\texttt{goto} змінною стану та великим \texttt{switch}.

\smallhead{2.4. Приведіть код до читабельного вигляду}

Застосуйте однакові відступи, змістовні імена та короткі коментарі лише там,
де причина дії неочевидна. Структура програми має читатися згори вниз.

\smallhead{Додаткове завдання • необов’язкове}

Винесіть один завершений етап алгоритму в окрему функцію та стисло поясніть,
як це покращило структуру без зміни результату.

\newpage

\eyebrow{Самоконтроль і робота з помилкою}
\pagetitle{3. Перевірка результату}

\smallhead{3.1. Перевірте відсутність goto}

Скористайтеся пошуком у файлі за словом \texttt{goto}. У фінальному коді не
повинно залишитися ані операторів \texttt{goto}, ані міток, призначених для
таких переходів. Переконайтеся, що цикли мають явні умови завершення.

\smallhead{3.2. Скомпілюйте структурну версію}

\cmd{gcc -std=c11 -Wall -Wextra -Wpedantic task\_XX\_structured.c -o structured}

Виправте всі помилки й попередження. Потім збережіть результат запуску:

\cmd{.\textbackslash structured.exe | Tee-Object structured\_output.txt}

Для Linux або macOS використайте
\texttt{./structured | tee structured\_output.txt}.

\smallhead{3.3. Порівняйте виведення}

У Windows PowerShell виконайте:

\cmd{Compare-Object (Get-Content original\_output.txt) (Get-Content structured\_output.txt)}

У Linux або macOS виконайте:

\cmd{diff -u original\_output.txt structured\_output.txt}

За повного збігу обидві команди не виводять відмінностей. Сам факт успішної
компіляції не доводить збереження поведінки.

\smallhead{3.4. Дослідіть одну помилку}

У тимчасовій копії структурного файлу змініть умову завершення одного циклу
так, щоб виникла зайва або пропущена ітерація. Скомпілюйте, порівняйте
виведення, поясніть відмінність і видаліть тимчасову копію. Фінальний файл має
залишитися правильним.

\subhead{Запитання для самоконтролю}

\begin{enumerate}
  \item Які три базові керівні структури використовує структурне програмування?
  \item Чому велика кількість довільних переходів ускладнює аналіз програми?
  \item Як розпізнати цикл, організований переходом назад?
  \item Коли під час перероблення доречні \texttt{break} і \texttt{continue}?
  \item Чим \texttt{while} відрізняється від \texttt{do--while}?
  \item Що означає поведінкова еквівалентність двох програм?
  \item Чому одного тестового набору може бути недостатньо для доведення
        еквівалентності?
\end{enumerate}

\newpage

\eyebrow{Обов’язкові розділи}
\pagetitle{4. Структура звіту}

\subhead{1. Титульний аркуш}

Оформіть за зразком на с. 7. Замініть поля у квадратних дужках власними
даними. Назву дисципліни та тему лабораторної роботи не скорочуйте.

\subhead{2. Тема та мета роботи}

Укажіть номер і тему лабораторної роботи. Сформулюйте мету через аналіз
неструктурованого керування, застосування структурних конструкцій і перевірку
еквівалентності.

\subhead{3. Варіант і середовище}

Укажіть номер та повне ім’я призначеного файлу, операційну систему, редактор і
компілятор із версією. Версію отримайте командою \texttt{gcc --version} або
\texttt{clang --version}.

\subhead{4. Аналіз початкової програми}

Стисло опишіть призначення програми, початкові дані та результат. Додайте
схему або таблицю переходів. Поясніть, які групи міток реалізують цикли,
розгалуження, пропуски та виходи.

\subhead{5. План перетворення}

Для кожного суттєвого фрагмента зазначте структурну заміну. Зручно подати
таблицю:

\renewcommand{\arraystretch}{1.25}
\begin{tabularx}{\textwidth}{|>{\RaggedRight\arraybackslash}p{0.24\textwidth}|>{\RaggedRight\arraybackslash}p{0.29\textwidth}|X|}
\hline
\textbf{Початковий фрагмент} & \textbf{Роль переходу} & \textbf{Структурна заміна} \\
\hline
Назви міток і переходів & Що саме реалізовано & Обрана конструкція та пояснення \\
\hline
\end{tabularx}

\subhead{6. Структурований код}

Наведіть повний фінальний текст \filename{task\_XX\_structured.c}. Код у звіті
має точно відповідати електронному файлу. Початковий код повторно друкувати не
потрібно: достатньо його аналізу та окремого електронного файла.

\subhead{7. Перевірка, відповіді та висновки}

Наведіть команди компіляції, обидва тексти виведення та результат порівняння.
Дайте короткі відповіді на запитання зі с. 4. У висновках зазначте, які
переходи було найскладніше перетворити і як перевірено збереження поведінки.

\newpage

\eyebrow{Єдині вимоги для цієї роботи}
\pagetitle{5. Вимоги до оформлення}

\subhead{Формат документа}

Звіт українською мовою. Формат А4, книжкова орієнтація. Поля: ліве 25 мм,
праве 15 мм, верхнє й нижнє 20 мм. Основний текст: Times New Roman, 14 пт,
міжрядковий інтервал 1,5; абзацний відступ 1,25 см.

\subhead{Заголовки та сторінки}

Заголовки розділів — напівжирні, 14 пт. Титульний аркуш — окрема перша
сторінка без видимого номера. Решту сторінок нумеруйте внизу, починаючи з 2.
Фіксований обсяг не встановлюється: усі обов’язкові розділи мають бути
змістовними й повними.

\subhead{Схема або таблиця переходів}

Схема має бути читабельною після друку. Для кожної мітки покажіть вхідні
переходи, умови та подальший напрям виконання. Якщо використовуєте таблицю,
не обмежуйтеся переліком назв: поясніть роль кожного переходу в алгоритмі.

\subhead{Код програми}

Код подавайте текстом моноширинним шрифтом Consolas або Courier New,
10–12 пт, з одинарним інтервалом і збереженими відступами. Не замінюйте код
знімками екрана. Довгі рядки переносьте так, щоб вони не виходили за поля.

\subhead{Результати запуску}

Вставте в звіт вміст \filename{original\_output.txt} і
\filename{structured\_output.txt}, а також повідомлення засобу порівняння.
Для кожного блоку зазначте, якою командою його отримано. Виведення подавайте
моноширинним шрифтом 10–12 пт.

\subhead{Доказ самостійного аналізу}

У звіті мають бути власні пояснення відповідності між початковими переходами
та структурними конструкціями. На захисті студент повинен показати у своєму
варіанті принаймні один цикл, одне розгалуження та один вихід або пропуск,
після чого пояснити їх перетворення.

\subhead{Фінальний формат}

Підготуйте звіт у PDF для електронного подання. Якщо викладач вимагає
паперовий примірник, роздрукуйте його на А4, перевірте читабельність коду,
порядок сторінок і заповнення титульного аркуша.

Не додавайте до звіту сторонні варіанти, згенеровані відповіді без пояснення
або фрагменти, яких немає у вашій фінальній програмі.

\newpage

\thispagestyle{empty}
\begin{center}
{\fontsize{10}{13}\selectfont
МІНІСТЕРСТВО ОСВІТИ І НАУКИ УКРАЇНИ\par
\vspace{8pt}
Київський університет інтелектуальної власності та права\par
\vspace{8pt}
Кафедра кібербезпеки та інформаційних технологій\par}

\vspace{38mm}

{\fontsize{16}{19}\selectfont\bfseries ЗВІТ\par}
\vspace{4pt}
{\fontsize{13}{16}\selectfont\bfseries
про виконання лабораторної роботи № 1\par}

\vspace{13mm}

з дисципліни\par
{\bfseries «Інженерія програмного забезпечення\\
та мови програмування»\par}

\vspace{10mm}

Тема: «Імперативна та структурна\\
парадигми програмування»

\vspace{20mm}

\begin{minipage}{0.47\textwidth}
\raggedright
Виконав(ла):\par
студент(ка) 1 курсу, групи [шифр групи]\par
[Прізвище, ім’я, по батькові]\par
Варіант № [номер]\par

\vspace{10pt}
Перевірив:\par
Петров О.В., к.т.н., доцент\par
кафедри кібербезпеки та\par
інформаційних технологій
\end{minipage}

\vfill

Київ, 2026
\end{center}

\newpage

\eyebrow{Підсумкова перевірка}
\pagetitle{6. Здавання роботи}

\subhead{Склад комплекту}

У папці \filename{Lab01\_Прізвище\_Група} мають бути:

\begin{itemize}
  \item \filename{task\_XX\_original.c} — незмінний файл викладача;
  \item \filename{task\_XX\_structured.c} — структурована версія без
        \texttt{goto};
  \item \filename{original\_output.txt} і
        \filename{structured\_output.txt};
  \item \filename{Lab01\_Прізвище\_Група.pdf} — оформлений звіт.
\end{itemize}

Перед надсиланням заархівуйте папку у ZIP, якщо цього вимагає обраний
викладачем канал здавання. Будьте готові скомпілювати обидві програми та
пояснити перетворення без читання готового тексту.

\subhead{Критерії перевірки}

\textbf{Повнота:} виконано аналіз переходів, перероблення, компіляцію,
порівняння та всі розділи звіту.

\textbf{Правильність:} структурна програма зберігає обчислення, гілки,
кількість і порядок ітерацій та виводить той самий результат.

\textbf{Структурність:} у коді немає \texttt{goto}, прихованого програмного
лічильника або великого \texttt{switch}, що відтворює початкові мітки.

\textbf{Якість коду:} конструкції вибрано відповідно до змісту алгоритму,
відступи послідовні, імена та коментарі допомагають читанню.

\textbf{Звіт і захист:} електронні файли відповідають звіту, студент розуміє
свій варіант і може пояснити основні рішення.

\subhead{Коли потрібне доопрацювання}

Програма не компілюється; залишилися оператори \texttt{goto}; результат або
порядок повідомлень відрізняється; частину гілок вилучено; відповідь
підставлено константою; структурне керування замінено машиною станів;
початковий файл змінено; у звіті немає аналізу чи доказу порівняння.

\subhead{Перевірте перед здаванням}

Номер варіанта та назви файлів правильні. Обидві програми компілюються з
прапорцями попереджень. Виведення збігається. Пошук за словом \texttt{goto} у
структурному файлі не дає результатів. Титульний аркуш заповнено. У ZIP немає
виконуваних файлів, тимчасових копій або чужих варіантів.

Додаткове завдання не замінює обов’язкових пунктів. На захисті покажіть один
цикл і одне розгалуження до та після перетворення.

\newpage

\eyebrow{Додаток • Windows}
\pagetitle{7. Установлення GCC}

\smallhead{1. Установіть MSYS2}

Відкрийте офіційну сторінку \url{https://www.msys2.org/}, завантажте
інсталятор і виконайте стандартне встановлення. Після завершення відкрийте
середовище \textbf{MSYS2 UCRT64}.

\smallhead{2. Оновіть систему пакетів}

У терміналі MSYS2 UCRT64 виконайте:

\cmd{pacman -Syu}

Якщо термінал попросить завершити роботу, закрийте його, знову відкрийте
MSYS2 UCRT64 і повторіть команду.

\smallhead{3. Установіть набір інструментів}

\cmd{pacman -S --needed base-devel mingw-w64-ucrt-x86\_64-toolchain}

Підтвердьте запропонований набір пакетів. Він містить GCC та інші засоби,
необхідні для компіляції й налагодження.

\smallhead{4. Додайте GCC до PATH}

Додайте каталог \filename{C:\textbackslash msys64\textbackslash
ucrt64\textbackslash bin} до системної або користувацької змінної
\texttt{Path}. Повністю закрийте й знову відкрийте VS Code. У новому терміналі
перевірте:

\cmd{gcc --version}

\smallhead{5. Налаштуйте VS Code}

У розділі Extensions установіть розширення \textbf{C/C++} від Microsoft.
Редактор і розширення не містять власного компілятора: команда \texttt{gcc}
має працювати в терміналі незалежно від кнопки запуску VS Code.

\smallhead{Якщо компіляція не вдається}

Якщо \texttt{gcc} не розпізнається, перевірте шлях
\filename{C:\textbackslash msys64\textbackslash ucrt64\textbackslash bin} і
перезапустіть VS Code. Якщо компілятор не знаходить файл, перевірте поточну
папку командами \texttt{pwd} і \texttt{ls} у MSYS2 або \texttt{Get-Location}
і \texttt{Get-ChildItem} у PowerShell.

\smallhead{Офіційні джерела}

MSYS2: \url{https://www.msys2.org/docs/installer/}\par
Visual Studio Code: \url{https://code.visualstudio.com/docs/cpp/config-mingw}

\end{document}
